% !TEX root = ../main.tex
\chapter{Thông tin tái lập và danh mục kỹ thuật}

\section{Quy trình tái lập kiểm tra}

Các lệnh sau được thực hiện từ thư mục gốc mã nguồn DA2. \texttt{npm ci} phù hợp khi cần cài đúng dependency từ lockfile; lần đánh giá ngày 09/08/2026 sử dụng dependency đã cài và chạy trực tiếp \texttt{npm run check}. Kết quả ghi nhận là ESLint không có lỗi/cảnh báo, 41/41 test đạt và Vite 7.3.6 build 1.906 module thành công.

\begin{verbatim}
npm ci
npm run check
\end{verbatim}

Script \texttt{check} thực hiện tuần tự ba lệnh:

\begin{verbatim}
npm run lint
npm test
npm run build
\end{verbatim}

Để chạy ứng dụng phát triển, cần MongoDB và file biến môi trường hợp lệ. Mã nguồn cung cấp script chạy riêng frontend/backend và script điều phối:

\begin{verbatim}
npm run dev
npm start
npm run start:all
\end{verbatim}

Trước khi lưu số liệu kiểm thử chính thức, cần ghi lại ngày giờ, hệ điều hành, phiên bản Node/npm, commit/tag, trạng thái working tree và toàn bộ output của lệnh \texttt{npm run check}. Không đưa file \texttt{.env}, token, URI có credential hoặc dữ liệu CV thật vào phụ lục.

\section{Biến môi trường}

\begin{table}[H]
    \centering
    \small
    \begin{tabularx}{\textwidth}{|p{0.34\textwidth}|Y|} \hline
        \textbf{Tên biến} & \textbf{Mục đích} \\ \hline
        \texttt{PORT} & Cổng lắng nghe của backend. \\ \hline
        \texttt{MONGODB_URI}, \texttt{MONGO_URI} & Chuỗi kết nối MongoDB; \texttt{MONGO\_URI} được ưu tiên nếu cả hai cùng tồn tại. Giá trị có thể chứa thông tin xác thực và phải được bảo vệ. \\ \hline
        \texttt{MONGO_DB_NAME} & Tên cơ sở dữ liệu; mặc định là \texttt{RabbitCV}. \\ \hline
        \texttt{JWT_SECRET} & Khóa ký/xác minh JWT; production phải dùng giá trị ngẫu nhiên mạnh. \\ \hline
        \texttt{NODE_ENV} & Phân biệt môi trường phát triển và production. \\ \hline
        \texttt{CLIENT_URL} & Bổ sung nguồn frontend được CORS cho phép. \\ \hline
        \texttt{VITE_SOCKET_URL} & Ghi đè địa chỉ Socket.IO ở frontend khi không dùng địa chỉ mặc định. \\ \hline
        \texttt{AI_ENABLED} & Bật hoặc tắt toàn bộ nhóm chức năng AI. \\ \hline
        \texttt{AI_DAILY_LIMIT} & Quota AI theo user trong một ngày. \\ \hline
        \texttt{AI_REQUEST_TIMEOUT_MS} & Thời gian tối đa chờ provider cho một request. \\ \hline
        \texttt{GROQ_API_KEY} & Khóa truy cập Groq; chỉ đặt ở backend/secret store. \\ \hline
        \texttt{GROQ_MODEL} & Tên model dùng cho các workflow AI. \\ \hline
        \texttt{CRON_SECRET} & Khóa Bearer bảo vệ endpoint Vercel Cron ở production. \\ \hline
    \end{tabularx}
    \caption{Danh sách biến môi trường không kèm giá trị bí mật}
    \label{tab:environment-variables}
\end{table}

\section{Danh mục API theo miền nghiệp vụ}

Bảng \ref{tab:api-catalog} và \ref{tab:api-catalog-part2} dùng prefix để tránh lặp đường dẫn. Ký hiệu \texttt{:id}, \texttt{:jobId} và \texttt{:room} là tham số động.

\begin{table}[H]
    \centering
    \footnotesize
    \begin{tabularx}{\textwidth}{|p{0.23\textwidth}|p{0.30\textwidth}|Y|} \hline
        \textbf{Miền} & \textbf{Prefix/tài nguyên} & \textbf{Phương thức và thao tác chính} \\ \hline
        Vận hành & \texttt{/api/health}, \texttt{/api/admin/health/details}, \texttt{/api/cron/job-expiry} & GET trạng thái tối giản; admin đọc chẩn đoán; Vercel Cron gọi kiểm tra hạn tin bằng Bearer secret. \\ \hline
        Xác thực & \texttt{/api/register}, \texttt{/api/login}, \texttt{/api/logout}, \texttt{/api/me} & POST đăng ký, đăng nhập, đăng xuất; GET phiên hiện tại. \\ \hline
        Hồ sơ User & \texttt{/api/user}, \texttt{/api/users} & GET tìm kiếm/xem profile; PUT profile, security và password. \\ \hline
        Resume & \texttt{/api/resume}, \texttt{/api/resumes} & GET danh sách/chi tiết/view; POST tạo hoặc cập nhật; DELETE theo id. \\ \hline
        Job & \texttt{/api/jobs} & GET danh sách, chi tiết và Job doanh nghiệp; POST tạo; PUT sửa; DELETE theo id. \\ \hline
        Application & \texttt{/api/jobs/apply}, \texttt{/api/jobs/applications} & POST nộp Resume/file; GET hồ sơ; PUT trạng thái Application. \\ \hline
        Saved jobs & \texttt{/api/saved-jobs} & GET danh sách; POST lưu; DELETE bỏ lưu theo jobId. \\ \hline
    \end{tabularx}
    \caption{Danh mục API theo miền nghiệp vụ (phần 1)}
    \label{tab:api-catalog}
\end{table}

\begin{table}[H]
    \centering
    \footnotesize
    \begin{tabularx}{\textwidth}{|p{0.23\textwidth}|p{0.30\textwidth}|Y|} \hline
        \textbf{Miền} & \textbf{Prefix/tài nguyên} & \textbf{Phương thức và thao tác chính} \\ \hline
        Message/chat & \texttt{/api/messages}, \texttt{/api/chat} & GET/POST message; GET user/thread/unread; PATCH đọc hoặc ẩn room. \\ \hline
        Notification & \texttt{/api/notifications} & GET danh sách/số chưa đọc; PUT đọc từng mục/tất cả; DELETE theo id. \\ \hline
        AI & \texttt{/api/ai} & GET trạng thái/model/quota; POST viết tiểu sử, review, điền CV, chấm CV--Job và từng lượt phỏng vấn. \\ \hline
        Admin Job & \texttt{/api/admin/jobs} & GET danh sách; PUT hide/unhide; DELETE theo id. \\ \hline
        Admin Company & \texttt{/api/admin/companies} & GET danh sách; PUT trạng thái verification. \\ \hline
        Admin User & \texttt{/api/admin/users} & PUT ban; DELETE theo id. \\ \hline
        Admin thống kê & \texttt{/api/admin/stats} & GET số liệu dashboard. \\ \hline
    \end{tabularx}
    \caption{Danh mục API theo miền nghiệp vụ (phần 2)}
    \label{tab:api-catalog-part2}
\end{table}

\section{Bất biến dữ liệu cần bảo vệ}

\begin{table}[H]
    \centering
    \small
    \begin{tabularx}{\textwidth}{|p{0.28\textwidth}|Y|} \hline
        \textbf{Bất biến} & \textbf{Ý nghĩa kiểm thử} \\ \hline
        Application duy nhất theo Job--candidate & Hai request gần nhau vẫn không tạo hai hồ sơ cho cùng ứng viên và Job. \\ \hline
        Application có \texttt{jobSnapshot} & Lịch sử còn tiêu đề, công ty và ngữ cảnh cần thiết sau khi Job bị xóa. \\ \hline
        Resume ảo không thuộc thư viện thường & PDF nộp cho một Application không làm rác danh sách CV quản lý của ứng viên. \\ \hline
        Danh sách Resume không trả \texttt{pdfData} & Truy vấn metadata không kéo theo toàn bộ base64 và không lộ file ngoài nhu cầu. \\ \hline
        Trạng thái chat gắn với user--room & Đánh dấu đọc hoặc ẩn của một phía không thay đổi dữ liệu cá nhân của phía còn lại. \\ \hline
        AiUsage không cần nguyên prompt & Quota và quan sát kỹ thuật dùng metadata/hash; dữ liệu CV đầy đủ không được lưu lại chỉ vì mục đích thống kê. \\ \hline
        Rich text CV qua danh sách cho phép & Thẻ định dạng được giữ nhưng thuộc tính, script, iframe, object và nội dung nhúng thực thi phải bị loại ở server. \\ \hline
        Cron không gửi thông báo lặp theo trạng thái & Các mốc đã thông báo chỉ được cập nhật sau khi tạo notification thành công; cần test thêm gọi đồng thời. \\ \hline
    \end{tabularx}
    \caption{Các bất biến quan trọng cho integration test}
    \label{tab:data-invariants}
\end{table}

\section{Rubric đề xuất cho chất lượng AI}

Rubric này dành cho lần đánh giá tiếp theo, không phải số liệu đã thu được trong báo cáo. Mỗi output được chấm theo thang 1--5 rồi quy đổi theo trọng số.

\begin{table}[H]
    \centering
    \small
    \begin{tabularx}{\textwidth}{|p{0.28\textwidth}|p{0.15\textwidth}|Y|} \hline
        \textbf{Tiêu chí} & \textbf{Trọng số} & \textbf{Câu hỏi chấm} \\ \hline
        Bám sát dữ liệu nguồn & 25\% & Output có giữ đúng sự kiện trong CV/Job và tránh tự thêm tên, thời gian, bằng cấp hay thành tích không có nguồn không? \\ \hline
        Đúng đắn và nhất quán & 20\% & Nhận xét có mâu thuẫn với dữ liệu, với chính nó hoặc với ràng buộc workflow không? \\ \hline
        Khả năng hành động & 20\% & Người dùng có thể sửa CV hoặc luyện phỏng vấn dựa trên đề xuất cụ thể không? \\ \hline
        Rõ ràng và phù hợp & 15\% & Ngôn ngữ có dễ hiểu, đúng vai trò, đúng độ dài và đúng ngữ cảnh nghề nghiệp không? \\ \hline
        An toàn, riêng tư, công bằng & 15\% & Output có tiết lộ dữ liệu không cần thiết, suy đoán thuộc tính nhạy cảm hoặc tạo lời khuyên phân biệt đối xử không? \\ \hline
        Tuân thủ định dạng & 5\% & JSON/schema, trường bắt buộc, giới hạn danh sách và kiểu dữ liệu có đúng không? \\ \hline
    \end{tabularx}
    \caption{Rubric 100 điểm đề xuất cho output AI}
    \label{tab:ai-rubric}
\end{table}

Một quy trình tối thiểu nên dùng ít nhất 20 CV đã ẩn danh, trải đều sinh viên, junior và mid-level; ít nhất 10 Job description; hai người chấm độc lập; cùng một model, prompt và tham số. Ngoài điểm trung bình, cần báo cáo tỷ lệ output bị loại vì bịa đặt/vi phạm an toàn và mức đồng thuận giữa hai người chấm. Mẫu nào chứa dữ liệu cá nhân thật chỉ được sử dụng khi có quyền hợp lệ và kế hoạch xóa sau đánh giá.

\section{Checklist trước khi nộp phiên bản chính thức}

\begin{itemize}
    \item Commit toàn bộ mã nguồn DA2 dùng làm bằng chứng, loại tệp tạm và gắn tag phiên bản.
    \item Chạy lại \texttt{npm run check}; lưu output và cập nhật số liệu nếu thay đổi.
    \item Chạy checklist thủ công ở Bảng \ref{tab:manual-checklist} và lưu ảnh/kết quả thực tế.
    \item Xác minh tên đề tài, sinh viên, giảng viên, loại báo cáo và tháng nộp trên bìa.
    \item Biên dịch LaTeX từ đầu, kiểm tra citation/reference và xem trực quan toàn bộ PDF.
    \item Bắt buộc cấu hình JWT secret ở production; vô hiệu hóa việc tự tạo tài khoản quản trị trình diễn và không công bố mật khẩu trong Git/PDF.
    \item Kiểm thử quyền vai trò ở API AI, lưu việc, ứng tuyển; kiểm tra thành viên phòng và giới hạn ảnh chat ở server.
    \item Loại secret, dữ liệu CV thật, cookie, URI có thông tin xác thực và file môi trường khỏi Git/PDF.
\end{itemize}
